% =============================================================================
% Configuración general del documento
% =============================================================================

% Idioma, codificación y tipografía
\usepackage[utf8]{inputenc}
\usepackage[T1]{fontenc}
\usepackage[provide=*,spanish]{babel}
\AtBeginDocument{\ifdefined\decimalpoint\decimalpoint\fi}
\usepackage{lmodern}

% Página y composición tipográfica
\usepackage[
  a4paper,
  top=2.8cm,
  bottom=2.8cm,
  left=3.2cm,
  right=2.8cm,
  headheight=14.5pt
]{geometry}
\usepackage{microtype}
\usepackage{setspace}
\onehalfspacing
\setlength{\parindent}{1.25cm}
\setlength{\parskip}{0pt}
\setlength{\emergencystretch}{2em}
\raggedbottom

% Matemáticas y unidades
\usepackage{amsmath,amssymb,mathtools}

% Figuras y tablas
\usepackage{graphicx}
\graphicspath{{Figuras/}}
\usepackage{booktabs}
\usepackage{longtable}
\usepackage{tabularx}
\usepackage{array}
\usepackage{multirow}
\usepackage{caption}
\usepackage{subcaption}
\usepackage{float}
\usepackage{placeins}
\captionsetup{font=small,labelfont=bf}

% Listas y código fuente
\usepackage{enumitem}
\usepackage{listings}
\usepackage{xcolor}
\definecolor{codigoFondo}{RGB}{247,248,250}
\definecolor{codigoComentario}{RGB}{75,120,75}
\lstset{
  basicstyle=\ttfamily\small,
  backgroundcolor=\color{codigoFondo},
  commentstyle=\color{codigoComentario},
  frame=single,
  breaklines=true,
  showstringspaces=false,
  tabsize=4
}

% Títulos de capítulos y secciones
\usepackage{titlesec}
\titleformat{\chapter}[hang]
  {\normalfont\bfseries\Large}
  {\thechapter.}
  {0.75em}
  {}
\titleformat{name=\chapter,numberless}[hang]
  {\normalfont\bfseries\Large}
  {}
  {0pt}
  {}
\titlespacing*{\chapter}{0pt}{0pt}{1.6em}
\titlespacing*{\section}{0pt}{2.2ex plus 0.5ex minus 0.2ex}{1.0ex}
\titlespacing*{\subsection}{0pt}{2.0ex plus 0.5ex minus 0.2ex}{0.8ex}
\setcounter{secnumdepth}{3}
\setcounter{tocdepth}{3}

% Encabezados y pies de página
\usepackage{fancyhdr}
\fancypagestyle{plain}{%
  \fancyhf{}
  \fancyfoot[C]{\thepage}
  \renewcommand{\headrulewidth}{0pt}
  \renewcommand{\footrulewidth}{0pt}
}

% Citas y referencias. En Overleaf se emplea biblatex-apa con Biber.
% El segundo bloque permite compilar en instalaciones mínimas de LaTeX.
\usepackage{csquotes}
\IfFileExists{biblatex.sty}{%
  \usepackage[
    backend=biber,
    style=apa,
    sorting=nyt,
    sortcites=true,
    maxbibnames=99,
    giveninits=true,
    uniquename=false
  ]{biblatex}
  \DeclareLanguageMapping{spanish}{spanish-apa}
  \addbibresource{Figuras/referencias.bib}
  \newcommand{\imprimirreferencias}{%
    \cleardoublepage
    \printbibliography[heading=bibintoc,title={Referencias}]%
  }
}{%
  \usepackage[authoryear,round]{natbib}
  \bibliographystyle{apalike}
  \providecommand{\parencite}[1]{\citep{##1}}
  \providecommand{\textcite}[1]{\citet{##1}}
  \newcommand{\imprimirreferencias}{%
    \cleardoublepage
    \phantomsection
    \addcontentsline{toc}{chapter}{Referencias}
    \bibliography{referencias}%
  }
}


% Soporte de las figuras SVG y los algoritmos de la sección 3.1.
\usepackage{svg}
\svgpath{{Figuras/}}
\newcommand{\FiguraMetodologia}[2]{%
  \includegraphics[width=#1]{Figuras/#2.pdf}%
}
\usepackage{algorithm}
\usepackage{algpseudocode}
\numberwithin{algorithm}{chapter}
\floatname{algorithm}{Algoritmo}
\algrenewcommand\algorithmicrequire{\textbf{Entrada:}}
\algrenewcommand\algorithmicensure{\textbf{Salida:}}
\algrenewcommand\algorithmicif{\textbf{si}}
\algrenewcommand\algorithmicthen{\textbf{entonces}}
\algrenewcommand\algorithmicelse{\textbf{en otro caso}}
\algrenewcommand\algorithmicfor{\textbf{para}}
\algrenewcommand\algorithmicdo{\textbf{hacer}}
\algrenewcommand\algorithmicend{\textbf{fin}}
\algrenewcommand\algorithmicreturn{\textbf{devolver}}
\algrenewcommand\algorithmicfunction{\textbf{función}}

% Hipervínculos y referencias cruzadas: deben cargarse al final.
\usepackage[hidelinks]{hyperref}
\Urlmuskip=0mu plus 1mu\relax
\usepackage{bookmark}
\usepackage[capitalise,noabbrev]{cleveref}
\hypersetup{
  pdftitle={Estudio reproducible de escalabilidad por resolución en clasificación tridimensional con representación basada en octrees: enfoque clásico y red profunda},
  pdfauthor={Ricardo Andrés López Tarazona y Daniel Tarazona Sánchez},
  pdfsubject={Trabajo de grado - Ingeniería de Sistemas},
  pdfkeywords={clasificación tridimensional, octree, Net5, ModelNet40, reproducibilidad}
}

% Nombres en español
\addto\captionsspanish{%
  \renewcommand{\contentsname}{Índice general}
  \renewcommand{\tablename}{Tabla}
  \renewcommand{\listtablename}{Índice de tablas}
  \renewcommand{\listfigurename}{Índice de figuras}
  \renewcommand{\bibname}{Referencias}
}
\crefname{figure}{figura}{figuras}
\Crefname{figure}{Figura}{Figuras}
\crefname{table}{tabla}{tablas}
\Crefname{table}{Tabla}{Tablas}
\crefname{section}{sección}{secciones}
\Crefname{section}{Sección}{Secciones}
\crefname{chapter}{capítulo}{capítulos}
\Crefname{chapter}{Capítulo}{Capítulos}
\crefname{equation}{ecuación}{ecuaciones}
\Crefname{equation}{Ecuación}{Ecuaciones}

% Comandos útiles del proyecto
\newcommand{\resolucion}[1]{\ensuremath{#1^{3}}}
\newcommand{\ModelNet}{\textit{ModelNet40}}
\newcommand{\NetCinco}{\textit{Net5}}
\newcommand{\ContenidoPendiente}{%
  \par\noindent
  {\small\color{gray}\textit{Contenido pendiente de incorporación y revisión.}}%
  \par\vspace{0.45\baselineskip}
}

% El logotipo PDF tiene prioridad. El JPG permite que este paquete compile
% inmediatamente y puede eliminarse cuando se copie Figuras/logo_uis.pdf.
\newcommand{\LogoUIS}[1][0.23\textwidth]{%
  \IfFileExists{Figuras/logo_uis.pdf}{%
    \includegraphics[width=#1]{Figuras/logo_uis.pdf}%
  }{%
    \IfFileExists{Figuras/logo_uis.jpg}{%
      \includegraphics[width=#1]{Figuras/logo_uis.jpg}%
    }{%
      \fbox{\parbox[c][2.1cm][c]{4.2cm}{\centering Logo UIS\\pendiente}}%
    }%
  }%
}

\makeatletter\def\theHALG@line{\thealgorithm.\arabic{ALG@line}}\makeatother
